\section{Conclusion}
% Recap findings, future work (robustness, multilingual, watermarking)
% In this work, we constructed the UNLamb benchmark of 11.6k question–answer pairs to test how fact popularity affects unlearning in LLMs. Our experiments across four state-of-the-art unlearning methods, three pairs of popular/rare forget sets, and three model sizes show that popular facts are significantly harder to remove and more likely to induce catastrophic forgetting, whereas rare facts can be erased more cleanly. These results validate that fact popularity matters, expose limitations of current algorithms, and provide a foundation for adaptive schemes—e.g., dynamic learning-rate selection or popularity-aware, scale-sensitive algorithm choice—to make unlearning more precise and stable.
In this work, we constructed the UNLamb benchmark of 11.6k QA pairs to study how fact popularity affects unlearning in LLMs. Across four methods, three popular/rare forget-set pairs, and three model sizes, we find popular facts harder to remove and likely to cause catastrophic forgetting, while rare facts are erased more cleanly. These results show that popularity matters, expose algorithmic limitations, and lay the foundations for adaptive schemes such as dynamic learning rate selection and popularity-aware method choice to improve unlearning precision and stability.





\section{Limitations}
Our findings reveal a critical limitation in current unlearning methods: they are blind to fact popularity. This work establishes a foundation to address this gap by introducing UNLamb, a benchmark designed to catalyze research into a new class of popularity-aware unlearning algorithms. Future work should prioritize adaptive strategies, such as developing dynamic learning rates or methods that select algorithms based on a fact's entrenchment. The development of such techniques is crucial for enabling the safe, precise, and reliable knowledge removal required for deploying unlearning in real-world LLMs.

